% 本文件由 LaTeX 排版 Agent 根据有效 translation.json 自动生成。
% author=Jerome; work_id=3007; title=荣福玛利亚的卒世童贞

\chapter{荣福玛利亚的卒世童贞}

% structure: p0001 source=h1 target=omitted reason=page-title-already-rendered-by-parent

\Emph{驳斥赫尔维迪乌斯}。\par

这篇短论约在公元三八三年问世。引发争论的问题是：我等主的母亲在祂诞生后是否仍保持童贞。赫尔维迪乌斯主张，福音中提及我等主的\Quote{姊妹}和\Quote{兄弟}，证明荣福童贞后来生育了子女，并以戴尔都良和维克托利努斯的著述支持其观点。其看法之结果是：童贞的地位被置于婚姻之下。热罗尼莫猛烈地站在另一边，试图证明所称的\Quote{姊妹}和\Quote{兄弟}，或为若瑟前妻所生之子女，或为表亲，即童贞之姊妹的子女。关于这场争辩的详细记述，见法勒\WorkTitle{基督教的早期岁月}第124页及以下。热罗尼莫撰写此篇时，他和赫尔维迪乌斯都在罗马，达玛稣任教宗。赫尔维迪乌斯唯一留存的当代记录，即热罗尼莫在以下篇章中的记载。\par

热罗尼莫驳斥赫尔维迪乌斯时提出三项主张：——\par

第一，若瑟只是名义上、而非真正的玛利亚的丈夫。\par

第二，主的\Quote{兄弟}是祂的表亲，而非亲兄弟。\par

第三，童贞优于婚姻状态。\par

第一项主张占据第三章至第八章。其依据是玛窦福音第一章十八至二十五节的记载，特别是\Quote{他们未曾同居}（第四章）、\Quote{若瑟…没有认识她，直到…}（第五章至第八章）等词句。\BibleRef{玛 1:18-25}\par

第二项主张占据第九章至第十七章，主要围绕\Quote{头胎儿子}（第九章、第十章）一词展开。热罗尼莫论证道，此词不仅适用于多子中的长子，也适用于独子；并论及所称的兄弟和姊妹，热罗尼莫断言他们是克娄帕（或克罗帕）的妻子玛利亚的子女（第十一章至第十六章）。他援引许多教会作家支持此观点（第十七章）。\par

第三项主张中，热罗尼莫为支持童贞优于婚姻的观点，论证道，不仅玛利亚、连若瑟也保持了童贞状态（第十九章）；婚姻虽有时是神圣的，却对祈祷构成很大障碍（第二十章），而圣经的教训是：童贞和节制之身的状态较婚姻更符合天主的旨意（第二十一章、第二十二章）。\par

不久前，某些弟兄要求我答复赫尔维迪乌斯所写的一本小册子。我一直拖延未作，并非因为维护真理并驳斥一个几乎不曾知道学问的第一线曙光、粗鲁无知的村夫是件难事，而是因为我担心我的答复会让他显得值得击败。此外还有一层考虑：一个骚动不安的家伙，世界上唯一自以为是司铎又在俗的人（如常言道，他以为口若悬河在于喋喋不休，并认为说某人坏话是良心无亏的见证），如果给他机会争辩，他会比以前更厉害地亵渎。他仿佛站在高台上，把他的观点广为传播。也有理由担心，当真理抛弃他时，他会用辱骂的武器攻击对手。但这些沉默的动机虽然正当，却因弟兄们被他的狂言所厌恶而受的刺激，更公正地不再影响我。因此，福音的斧头必须放在那棵不结果实的树根上，把它和不结果实的叶子都扔进火里，以便从未学会说话的赫尔维迪乌斯最终学会保持沉默。\par

我必须呼求圣神用我的口表达祂的旨意，并捍卫荣福玛利亚的童贞。我必须呼求主耶稣保护祂曾居住十个月的子宫这神圣住所免受性交的猜疑。我还必须恳求天主圣父显示，祂圣子的母亲——在成为新娘之前已是母亲——在生子之后仍继续是童贞。我们无意在口才的田野上驰骋，也不诉诸逻辑家的陷阱或亚里士多德的密林。我们将引出圣经的实际话语。让他被他自己用来攻击我们的同一证据所驳倒，这样他就能明白，他可能读过所写的内容，却无法辨别健全信仰的确立结论。\par

他的第一个陈述是：\Quote{玛窦写道：\BibleQuote{耶稣基督的诞生是这样的：祂的母亲玛利亚已许配给若瑟，他们尚未同居，她因圣神怀了孕。她的丈夫若瑟是个义人，不愿公开羞辱她，有意悄悄休了她。但当他在思索这些事时，看，上主的天使在梦中显现给他，说：若瑟，达味之子，不要怕娶你的妻子玛利亚，因为那在她内受孕的是出于圣神。}}请注意，他说，所用的词是\Emph{许配}，而不是像你所说的\Emph{受托}，当然她之所以被许配，唯一的原因是她有一天要结婚。并且福音作者不会说\Quote{他们尚未同居}，如果他们将来不打算同居的话，因为没有人会对一个不打算吃饭的人用\Quote{他尚未吃饭}的说法。而且，天使称她为\Quote{妻子}，并说她与若瑟相结合。接着，我们被邀请聆听圣经的宣告：\BibleRef{玛 1:24}\BibleQuote{若瑟从睡梦中醒来，就遵从上主的天使所吩咐的，娶了他的妻子；但没有认识她，直到她生了儿子。}\par

4. 我们逐一审视这些论点，追踪这不敬之言的轨迹，以表明他自相矛盾。他承认她是聘妻，紧接着又称她为妻子——他既已承认她是聘妻，却又将她说成是某人的妻子。再者，他称她为妻子，然后又说她被聘的唯一原因是她终将成婚。而且，唯恐我们觉得这还不够，他说：\Quote{所用的词是\Emph{聘妻}，而非\Emph{受托}，也就是说，尚未成为妻子，尚未因婚姻的结合而联合。}但当他接着说：\Quote{福音作者绝不会将\Emph{尚未同居}一词用于那些不会同居的人，这正如一个人不会说\InnerQuote{他吃饭前}，而那人实际上并不打算吃饭一样。}我不知道是该悲伤还是该笑。我应该指摘他的无知，还是谴责他的鲁莽？这就好比，若有人说：\Quote{我在港口用餐前就航向了非洲}，除非他终有一天须在港口用餐，否则这句话就不能成立。如果我要说：\Quote{宗徒保禄在去西班牙之前被囚于罗马}，或者（我当然也可以说）\Quote{赫维迪乌在悔改之前就被死亡夺去}，那么保禄就必须获释后立即去西班牙，或者赫维迪乌必须在死后悔改，尽管圣经说：\Quote{在阴府中，谁能称谢你？}难道我们不应该理解，介词\Emph{\Quote{之前}}虽然常表示时间上的顺序，但有时仅指思想上的顺序吗？因此，如果有充分的原因阻止，我们的想法未必需要实现。因此，当福音作者说\Emph{\Quote{尚未同居}}时，他指的是结婚前的那一刻，表明事情已进展到如此地步，以至于那位聘妻即将成为妻子。仿佛他在说：在他们亲吻和拥抱之前，在婚姻完成之前，她被发现怀了孕。发现这一点的正是若瑟，他用丈夫般焦虑的目光——此时几乎是丈夫的特权——注视着聘妻隆起的小腹。然而，正如前面的例子所表明，这并不意味着他在玛利亚分娩后与她同房，因为他的欲望已因她已怀孕而熄灭。尽管我们在梦中被告知若瑟：\Quote{不要怕娶你的妻子玛利亚，}以及\Quote{若瑟从睡梦中醒来，就照主的天使所吩咐的办了，娶了他的妻子，}但无人应当因此感到困扰，仿佛因为称她为\Emph{妻子}，她就不再是\Emph{聘妻}，因为我们知道圣经中惯常将聘妻称为妻子。以下《申命记》的证据确立了这一点。\BibleRef{申 22:24}\BibleQuote{若有人在田野遇见已聘的少女，那人强与她同寝，他必须死，因为他玷辱了邻舍的妻子。}在另一处，\BibleRef{申 22:23}\BibleQuote{若有少女已是聘于人的处女，有人在城里遇见她，与她同寝，你们就要把这二人带到城门口，用石头砸死他们；少女因为她在城里没有喊叫，那人因为他玷辱了邻舍的妻子。这样，你就从你们中间除去这恶。}此外，\BibleRef{申 20:7}\BibleQuote{谁聘了妻，尚未迎娶呢？他可以回家去，免得他死在战场，别人娶了她。}若有人对童贞女为何在聘后而非在无人聘她（或用圣经的话说，没有丈夫）时怀孕有所怀疑，容我解释其原因有三：第一，藉着玛利亚的亲属若瑟的家谱，玛利亚的出身也可显明；第二，免得她按梅瑟法律被当作淫妇用石头砸死；第三，在她逃往埃及时，她可以有某种安慰，尽管是监护人多于丈夫。因为当时谁会相信童贞女的话，说她因圣神怀孕，且天使加俾额尔曾来宣告天主的旨意呢？岂不会所有人都像对待苏撒纳那样，视她为淫妇而加以定罪么？因为即使在今日，当全世界都已接受信仰时，犹太人仍然争论说，当依撒意亚说\Quote{看，一位童贞女将怀孕生子}时，希伯来词意指年轻女子，而非童贞女，即所用之词是\OriginalTerm{少女（希伯来语）}{Almah}，而非\OriginalTerm{童贞女（希伯来语）}{Bethulah}，这一点我们将在后面更详细地讨论。最后，除了若瑟、依撒伯尔、玛利亚本人，以及可能从他们那里听到真相的少数人之外，所有人都认为耶稣是若瑟的儿子。甚至福音作者们，在表达普遍看法（这是历史学家的正确规则）时，也称他为救主的父亲，例如\BibleRef{路 2:27}\BibleQuote{他（即西默盎）在圣神感动下进入圣殿；当父母抱着孩子耶稣进来，要按律法的惯例为他行礼时，}以及其他地方，\BibleRef{路 2:41}\BibleQuote{他的父母每年逾越节往耶路撒冷去。}随后：\Quote{过完了节，他们回去的时候，孩童耶稣仍留在耶路撒冷，他的父母并不知道。}也请注意玛利亚本人——她曾对加俾额尔回答说：\Quote{这事如何成就？因为我不认识男人。}——关于若瑟她说了什么：\BibleRef{路 2:48}\BibleQuote{孩子，你为什么这样对待我们？看，你的父亲和我，忧苦地找你。}我们在这里并不是像许多人主张的那样，听到的是犹太人或者讥诮者的话。福音作者称若瑟为父亲，玛利亚承认他是父亲。这并不是说（如我之前所说）若瑟真是救主的父亲，而是为了保全玛利亚的名誉，他被所有人视为他的父亲，尽管在他听到天使的劝诫之前，\BibleRef{玛 1:20}\BibleQuote{若瑟，达味之子，不要怕娶你的妻子玛利亚，因为那在她内受孕的是出于圣神。}他曾想暗暗地休退她，这表明他深知所怀的孩子并非他的。但我们已说得足够，更多地是为了传授教导而非反驳对手，以说明为何若瑟被称为我主的父亲，为何玛利亚被称为若瑟的妻子。这也同时回答了为何某些人被称为他的兄弟。\par

5. 不过，这一点将在后面适当的地方讨论。我们现在必须转向其他问题。现在要讨论的经文是：\\newline{}\\BibleQuote{若瑟从睡梦中醒来，就照上主的天使所嘱咐的办了，娶了他的妻子；若瑟没有认识她，直到她生了一个儿子，给他起名叫耶稣。}\\newline{}这里首先，我们的对手完全不必如此详尽地证明\\Emph{认识}一词指的是性交，而不是理智上的认识：好像有人否认这一点，或者任何有理智的人会想象\\Person{赫尔维迪乌斯}费力反驳的那种愚蠢。然后他会教我们，\\Emph{直到}这个副词暗示一个固定的、明确的时间，当那个时间满足时，他说那件先前没有发生的事就发生了，就像我们面前的例子：\\newline{}\\BibleQuote{却没有认识她，直到她生了儿子。}\\newline{}他说，很明显，她在生产之后就被认识了，而那种认识只是由于她生了儿子而被延迟。为了捍卫他的立场，他堆积经文，像蒙眼角斗士一样挥舞剑，喋喋不休，最后只伤到了自己。\par

6. 我们的回答简要如下——在圣经的语言中，\\Emph{认识}和\\Emph{直到}这两个词可能具有双重含义。关于前者，他自己给了我们一篇论文来表明它必须指性交，并且没有人怀疑它常常用于理智的认识，例如：\\newline{}\\BibleQuote{孩童耶稣留在耶路撒冷，他的父母并不知道。}\\newline{}现在我们必须证明，正如他在一种情况下遵循了圣经的用法一样，关于\\Emph{直到}这个词，他完全被同一圣经的权威驳倒了，圣经经常用这个词表示一个固定的时间（他自己告诉过我们），也经常表示没有限制的时间，例如当天主借先知的口对某些人说：\\newline{}\\BibleRef{依 46:4}\\newline{}\\BibleQuote{直到你们年老，我仍是一样。}\\newline{}难道当他们年老时，祂就不再是天主了吗？救主在福音中告诉宗徒们：\\newline{}\\BibleRef{玛 28:20}\\newline{}\\BibleQuote{我同你们天天在一起，直到今世的终结。}\\newline{}那么，在世界终结之后，主会离开他的门徒吗？在审判十二支派的时候，他们正坐在十二个宝座上，难道他们会失去主的陪伴吗？此外，保禄宗徒写信给格林多人说：\\newline{}\\BibleQuote{基督作为初果，然后那些属于基督的人，在他来临时。那时，结局来临，当他将王权交给天主父，当他毁灭一切率领者、一切掌权者和大能者。因为祂必须为王，直到把一切仇敌放在祂脚下。}\\newline{}即使这段经文涉及主的人性，我们也不否认这些话是对那位忍受十字架并随后被命令坐在右边的那一位说的。那么，他说\\BibleQuote{因为祂必须为王，直到把一切仇敌放在祂脚下}是什么意思呢？难道主只当他的仇敌开始被置于他脚下时才统治，一旦他们被置于脚下，他就停止统治吗？当然，当他的仇敌开始被置于他脚下时，他的统治才开始完满。达味在第四首登阶歌中也这样说：\\newline{}\\BibleQuote{看，仆人的眼目，怎样仰望主人的手，看，婢女的眼目，怎样注视主母的手，我们的眼目也同样仰望上主，直到他怜悯我们。}\\newline{}那么，先知会仰望天主直到获得怜悯，一旦获得怜悯，他就把眼睛转向地上吗？尽管他在别处说：\\newline{}\\BibleQuote{我的眼睛渴望你的救恩，渴望你正义的诺言。}\\newline{}我可以积累无数这样的用法，用一大串证据淹没我们攻击者的冗长言词；但是，我只再添加几个，让读者自己去发现类似的例子。\par

7. 天主的言语在创世纪中说：\\newline{}\\BibleQuote{他们便把所有外邦的神像和耳朵上的环子交给雅各伯；雅各伯把它们埋在舍根附近的橡树下，失落到今日。}\\newline{}同样，在申命纪结束时，\\newline{}\\BibleRef{申 34:5}\\newline{}\\BibleQuote{上主的仆人梅瑟死在那里，就在摩阿布地，照上主所说的。上主将他葬在摩阿布地，面对贝特培敖尔的山谷；但没有人知道他的坟墓，直到今日。}\\newline{}我们当然必须理解\\Emph{直到今日}是指该历史著作的写作时间，无论你倾向于认为梅瑟是五书的作者，还是厄斯德拉重新编辑了它们。任何一种情况我都不反对。现在的问题是，\\Emph{直到今日}这个词是否指书籍的出版或写作时间，如果是，那么现在许多年过去了，他必须证明，要么埋在橡树下的偶像被找到了，要么梅瑟的坟墓被发现了；因为他顽固地主张，只要\\Emph{直到}所指示的时间点尚未达到，就不会发生的事，一旦那个时间点达到，就开始发生。他最好留意圣经的惯用语，并与我们一起理解（他正是在这里陷入泥沼的）：有些事如果不明说可能会显得模棱两可，但它们被清楚地暗示了，而另一些事则留给我们的理性去推断。因为如果，当事件还在记忆犹新，还有见过梅瑟的人活着时，他的坟墓就能不为人所知，那么经过这么多时代之后，这种情况就更加可能了。同样，我们必须解释关于若瑟的叙述。福音作者指出了可能引起某种丑闻的情况，即玛利亚没有被她的丈夫认识，直到她生产，他这样做是为了让我们更加确信，若瑟在因异象的意义存有怀疑而有所保留的那一位，在生产之后也没有认识她。\par

8. 简言之，我想知道的是，为什么\Person{若瑟}直到她分娩之日都克制自己？\Person{赫尔维狄乌斯}自然会回答说，因为他听到了天使说：\Quote{她怀的孕是从圣神来的}\BibleRef{玛 1:20}。而反过来我们则反驳说，他确曾听到天使说：\Quote{\Person{若瑟}，\Person{达味}之子，不要怕娶你的妻子\Person{玛利亚}}\BibleRef{玛 1:20}。他之所以被禁止休妻，是为了不让他以为她是淫妇。那么，难道他真的被命令不得与妻子同房吗？难道不清楚这警告是为了不让他与她分离吗？而这位义人，他说，当听到天主子在她腹中时，还敢想接近她吗？好极了！那么，我们就要相信，同一个人，因如此相信一个梦而不敢触碰他的妻子，后来，当他从牧羊人那里得知，主的天使从天上来到他们面前说：\Quote{不要害怕！看，我给你们报告一个为全民族的大喜讯，今天在\Place{达味城}中，为你们诞生了一位救世者，他是主默西亚}；当天上的军旅与他一同高唱\Quote{天主受享光荣于高天，主爱的人在世享平安}；\BibleRef{路 2:14}当他看见义人\Person{西默盎}怀抱婴儿高呼：\Quote{主啊！现在可照你的话，让你的仆人平安去吧！因为我亲眼看见了你的救恩}；当他看见女先知\Person{亚纳}、贤士、星辰、\Person{黑落德}和天使们时；\Person{赫尔维狄乌斯}，我敢说，你要我们相信\Person{若瑟}既然深知这些惊人的奇事，还敢触碰天主的圣殿、圣神的居所、他主的母亲吗？\Person{玛利亚}至少\Quote{把这一切事默存在自己心中}。\BibleRef{路 2:33}你不能无耻地说\Person{若瑟}不知道这些事，因为路加告诉我们：\Quote{他的父亲和母亲就惊异他关于\Person{耶稣}所说的话}。\BibleRef{路 2:33}而你却以惊人的厚颜争辩说，希腊文抄本有误，尽管几乎所有希腊作家都在他们的著作中留下了这样的读法，不仅如此，许多拉丁作家也采用了同样的措辞。我们现在无需考虑抄本之间的差异，因为旧约和新约的全部记录此后都被译成了拉丁文，我们必须相信源头的泉水比支流更纯净。\par

9. \Person{赫尔维狄乌斯}会回答：\Quote{依我之见，你说的不过是琐碎之言。你的论证纯属浪费口舌，而讨论显示出更多的诡辩而非真理。为何圣经不能像关于\Person{塔玛尔}和\Person{犹大}那样，说\BibleRef{创 38:26}\InnerQuote{他就再没有认识她}？\Person{玛窦}难道找不到表达他意思的词吗？他说：\InnerQuote{他不认识她，直到她生了儿子}。那么，他在她分娩之后，确实认识了她，那个他在她分娩之前一直克制不认识的人。}\par

10. 如果你如此好辩，你自己的思想现在就会成为你的主人。你不允许在分娩和同房之间有任何时间间隔。你不能说：\Quote{若有妇人怀孕生男孩，她就不洁净七天，像在月经污秽的日子一样不洁净。第八天，要给婴孩割阳皮。她要在产血得洁净的三十三天内，不可摸圣物}等等。按你的说法，\Person{若瑟}必须立即接近她，并受到\Person{耶肋米亚}\BibleRef{耶 5:8}的责斥：\Quote{他们像发情的公马，各向邻舍的妻子嘶鸣}。否则，\Quote{他不认识她，直到她生了儿子}这句话如何能成立，如果他等到另一次洁净的时期过去，如果他的情欲必须再忍受四十天的漫长等待？母亲必须带着产污未净，哀哭的婴儿由接生婆照顾，而丈夫拥抱疲惫的妻子。这样，婚姻生活必定开始，好使福音作者不被指责说谎。但天主决不让我们这样想救主的母亲和一位义人。没有接生婆协助他的诞生；没有妇人的多事干预。她亲手用襁褓包裹他，既是母亲又是接生婆，\Quote{把他放在马槽里}，\BibleRef{路 2:7}据记载，\Quote{因为在客栈中没有地方给他们}。这句话一方面驳斥了伪经记载的狂言，因为\Person{玛利亚}亲自用襁褓包裹他；另一方面使\Person{赫尔维狄乌斯}的淫欲观念成为不可能，因为客栈里没有适合夫妇同房的地方。\par

11. 对于他提出的关于\Emph{他们同居以前}和\Emph{他不认识她，直到她生了儿子}这些词句的论点，现在已经给出了充分的答复。如果我的答复要遵循他论证的顺序，我现在必须继续讨论第三点。他坚持说\Person{玛利亚}生了其他儿子，并引用了经文：\Quote{\Person{若瑟}也\Place{达味城}上去登记，同他订了婚的\Person{玛利亚}一同去，那时她怀了孕。他们在那里的时候，她分娩的日期满了，便生了她的头胎儿子}。他由此试图证明，\OriginalTerm{头胎}{first-born}一词只适用于有兄弟的人，正如\OriginalTerm{独生子}{only begotten}是指父母的独子。\par

12. 我们的立场如下：每一个独生子都是长子，但并非每一个长子都是独生子。我们理解长子不仅指有弟妹者，也指没有比他先出生的。\BibleRef{户 18:15} \BibleQuote{一切，}主对亚郎说，\BibleQuote{那开胎的，无论是人是牲畜，凡献给主的，应归你；然而人首生的，你必须赎回来；不洁净牲畜首生的，你也要赎回来。} 天主圣言定义\Emph{首生}为一切开胎的。否则，如果只有那些有弟弟的才可称为长子，那么司铎就不能要求头胎，直到后继者出生，因为如果没有后来的生产，它可能只是独生子而非长子。\BibleRef{户 18:16} \BibleQuote{凡当赎的，从一个月以上的，你要按估价，按圣所的舍客勒，五舍客勒（一舍客勒是二十季拉）赎回来。但牛、绵羊、山羊的头胎你不可赎，因为它们是圣的。} 天主圣言迫使我将一切开胎的奉献给天主，如果是洁净牲畜的头胎；若是不洁净的，我必须赎回来，并把价值交给司铎。我可能会回答：你为什么把我限制在一个月的短时间里？你为什么谈论长子，而我不知道是否有弟弟随后？等第二个出生吧。除非第二个出生使先前那个成为长子，否则我对司铎一无所欠。难道那些字句不会对我呼喊，指证我的愚昧，并宣告\Quote{长子}是开胎者的称号，而不是限于有弟弟者吗？那么，以若翰为例：我们都同意他是独生子；我想知道他是否也是长子，是否完全受法律的约束。这毫无疑问。无论如何，圣经这样论及救主：\BibleQuote{及至他们按梅瑟法律满了取洁的日子，他们便带着孩子上耶路撒冷去献给主（就如主的法律上所记载的：凡开胎首生的男性，应称为祝圣于主的），并按照主的法律所说，或一对斑鸠，或两只雏鸽，献上祭献。} 如果这法律只关乎长子，而长子只有在有后继者时才存在，那么没有人应该受长子法律的约束，如果他不知道是否有后继者。但是，既然没有弟弟的人也受长子法律的约束，我们得出结论：被称为长子的，是开胎且没有比他先出生的人，而不是出生后有弟弟的人。梅瑟在出谷纪中写道：\BibleRef{出 12:29} \BibleQuote{到了半夜，主击杀了埃及地所有头胎，从坐宝座的法郎的长子，直到地牢中囚犯的长子，以及一切牲畜的头胎。} 告诉我，那些被毁灭者击杀的人，只是你们的长子，还是也包括独生子？如果只有有兄弟的人才称为长子，那么独生子就免于死亡。如果事实上独生子也被杀了，那就违反了所宣布的判决，因为独生子与长子一同死亡。你要么将独生子从惩罚中释放出来——那样你就变得可笑；要么，如果你承认他们也被杀了，那么我们就赢得了论点，尽管我们不必感谢你：即独生子也被称为长子。\par

13. 赫尔维狄乌斯的最后一个论断就是这一点，也是他在讨论首生者时所希望证明的，即福音书中提到了主的兄弟。例如，\BibleRef{玛 12:46}\BibleQuote{看，他的母亲和他的兄弟站在外面，想要和他说话。} 又一处，\BibleRef{若 2:12}\BibleQuote{这事以后，他下到葛法翁，他和他的母亲、他的兄弟们。} 又，\BibleRef{若 7:3}\BibleQuote{他的兄弟们便对他说：你离开这里，往犹太去罢，好使你的门徒也看见你所行的事。因为没有人暗中行事，却愿意公开显扬。你如果行这些事，就向世界显扬你自己罢。} 若望接着补充说，\BibleRef{若 7:5}\BibleQuote{原来连他的兄弟也不相信他。} 马尔谷和玛窦也说：\Quote{他来到自己的家乡，在会堂里教训他们，以致他们惊讶说：这人从哪里得了这样的智慧和异能？这不是木匠的儿子吗？他的母亲不是叫玛利亚吗？他的兄弟们不是雅各伯、若瑟、西满和犹达吗？他的姊妹们不是都在我们这里吗？} 路加在宗徒大事录中也记载，\BibleRef{宗 1:14}\BibleQuote{这些人同几个妇女及耶稣的母亲玛利亚，并他的兄弟，都同心合意地恒切祈祷。} 宗徒保禄也与他们一致，并证实他们的历史真实性：\Quote{我原是藉启示上去的，可是别的宗徒，除了伯多禄和主的兄弟雅各伯，我都没有看见。} 又在一处，\BibleRef{格前 9:4}\BibleQuote{难道我们没有权利吃喝吗？难道我们没有权利带着妻子往来，如同其余的宗徒、主的兄弟和刻法一样吗？} 并且，惟恐有人不接受犹太人的证据（因为正是从他们口中我们听到他兄弟的名字），但坚持说他的同乡们对兄弟的看法也犯了他们对父亲的看法一样的错误，赫尔维狄乌斯便发出尖锐的警告，嚷道：\Quote{同样的名字在福音的另一处被重复，同样的人在那里是主的兄弟和玛利亚的儿子。} 玛窦说：\Quote{有许多妇女在那里（无疑是在主的十字架旁）远远地观看，她们是从加里肋亚跟随耶稣来服事他的；其中有玛利亚玛达肋纳、雅各伯和若瑟的母亲玛利亚，以及载伯德儿子的母亲。} 马尔谷也说：\Quote{也有妇女远远地观看，其中有玛利亚玛达肋纳、小雅各伯和若瑟的母亲玛利亚，以及撒罗默。} 同一处稍后又说：\Quote{还有好些同他上耶路撒冷来的妇女。} 路加也提到，\BibleRef{路 24:10}\BibleQuote{现在有玛利亚玛达肋纳、约安纳、雅各伯的母亲玛利亚，以及同她们一起的其他妇女。}\par

14. 我反复重复同样的事情，是为了防止他提出一个虚假的问题，并叫嚷说我隐瞒了对他有利的段落，而且他的观点不是被圣经的证据，而是被规避的论证所撕碎。看，他说，雅各伯和若瑟是玛利亚的儿子，就是被犹太人称为兄弟的同一批人。看，玛利亚是小雅各伯和若瑟的母亲。并且雅各伯被称为小者，以区别于大雅各伯，即载伯德的儿子，正如马尔谷在别处所说：\Quote{玛利亚玛达肋纳和若瑟的母亲玛利亚看见了安放他的地方。安息日一过，他们就买了香料，要去傅抹他。} 而且，正如所料，他说：\Quote{如果我们认为当其他妇女关心耶稣的安葬时，他的母亲玛利亚却不在场；或者如果我们杜撰某种第二位玛利亚——更何况\WorkTitle{圣若望福音}作证，当主在十字架上将她，作为他的母亲、现时是寡妇，托付给若望照顾时，她是在场的——那么我们对玛利亚持有一种多么可怜和不敬的看法！或者我们必须认为福音作者们错得如此离谱，以致误导我们，竟称玛利亚为那些被犹太人称为耶稣兄弟之人的母亲？}\par

15. 何等的黑暗，何等疯狂、冲向自我毁灭的狂怒！你说主的母亲在十字架旁，你说因她的寡居和孤独状况，她曾被托付给门徒\Person{若望}；仿佛按你自己的说法，她没有四个儿子和许多女儿，可以用她们的陪伴来安慰自己？你还把\Quote{寡妇}这名称加给她，这名称在圣经中是找不到的。尽管你引用了福音中的所有事例，唯独\Person{若望}的话不合你意。你顺带提到她在十字架旁，以免显得你故意省略此事，却只字未提与她在一起的其他妇女。如果你是无知，我可以原谅你，但我看出你沉默是有原因的。\newline{} 那么，请允许我指出\Person{若望}所说的，\BibleRef{若 19:25} \BibleQuote{在耶稣十字架旁，站着他的母亲，和他母亲的姊妹，克罗帕的玛利亚，以及玛利亚玛达肋纳。} 没人怀疑有两位宗徒都叫\Person{雅各伯}，\Person{载伯德}的儿子\Person{雅各伯}，和\Person{阿耳斐}的儿子\Person{雅各伯}。你是否认为那比较不出名的小\Person{雅各伯}，在圣经中被称为\Person{玛利亚}的儿子（但并非我主之母\Person{玛利亚}），是宗徒呢？如果他是宗徒，他必是\Person{阿耳斐}的儿子，并且相信\Person{耶稣}，\BibleQuote{因为他的兄弟也不信他。} 如果他不是宗徒，而是第三个\Person{雅各伯}（我无法知道他是谁），他怎能被视为主的兄弟？而且，作为第三个，他怎能被称为\Emph{小}以区别于\Emph{大}，因为\Emph{大}和\Emph{小}是用来表示两者之间的关系，而不是三者之间。此外，请注意，主的兄弟是一位宗徒，因为\Person{保禄}说：\BibleRef{迦 1:18} \BibleQuote{过了三年，我上耶路撒冷去见刻法，和他一起住了十五天。但除了主的兄弟雅各伯，我没有看见别的宗徒。} 并且在同一书信中，\BibleRef{迦 2:9} \BibleQuote{他们一看到所赐给我的恩宠，那被认为是柱石的雅各伯、刻法和若望，}等等。为使你不至于以为这\Person{雅各伯}是\Person{载伯德}的儿子，你只要读\WorkTitle{宗徒大事录}，就会发现后者早已被\Person{黑落德}杀害了。唯一的结论是，那被描述为小\Person{雅各伯}母亲的\Person{玛利亚}，是\Person{阿耳斐}的妻子，也是主母亲\Person{玛利亚}的姊妹，就是被福音作者\Person{若望}称为\BibleQuote{克罗帕的玛利亚}的那一位，无论这称呼是出自她的父亲、亲属，或是其他原因。但如果你认为她们是两个人，因为别处我们读\BibleQuote{小雅各伯的母亲玛利亚}，而这里\BibleQuote{克罗帕的玛利亚}，你仍须知道，在圣经中，同一个人有不同名字是常事。\Person{辣古耳}，\Person{梅瑟}的岳父，也被称为\Person{耶特洛}。\Person{基德红}，毫无明显原因地突然变成\Person{耶卢巴耳}。犹大王\Person{乌齐雅}另有一个名字\Person{阿札黎雅}。\Place{大博尔山}被称为\Place{伊塔比陵}。同样，\Place{赫尔孟山}被\Person{腓尼基人}称为\Place{萨尼尔}，被\Person{阿摩黎人}称为\Place{撒尼尔}。同一地区有三个名字：在\WorkTitle{厄则克耳}中称为\Place{乃格布}、\Place{特曼}和\Place{达隆}。\Person{伯多禄}也被称为\Person{西满}和\Person{刻法}。另一个福音中的\Person{热诚者犹达}被称为\Person{达陡}。还有许多其他例子，读者可以从圣经各部分自己收集。\par

16. 现在，我们有了我要努力说明的解释，即主的母亲玛利亚的姐妹的儿子们，（这些人当初不信，后来却信了）如何能被称为主的兄弟。可能的情况是，其中一个兄弟立即相信了，而其他人直到很久以后才相信，并且那一位玛利亚是雅各伯和若瑟的母亲，即\Quote{克罗帕的玛利亚}，她与阿尔斐的妻子是同一个人，另一位是次雅各伯的母亲。无论如何，如果她（后者）曾是主的母亲，圣若望就会像在其他地方一样给她这个称号，并且不会通过称她为其他儿子的母亲而给人错误的印象。但在这个阶段，我不愿争论克罗帕的妻子玛利亚和雅各伯与若瑟的母亲玛利亚是否为不同的妇女，只要清楚明白雅各伯和若瑟的母亲玛利亚与主的母亲不是同一个人。赫来维丢说，那么，你怎么证明那些并非他的兄弟的人被称为主的兄弟？我将说明那是怎么回事。在圣经中，有四种兄弟：本性的、民族的、同族的、情谊的。本性的兄弟的例子是厄撒乌和雅各伯、十二位圣祖、安德肋和伯多禄、雅各伯和若望。至于民族，所有犹太人彼此称为兄弟，正如在申命纪中，\BibleRef{申 15:12} \BibleQuote{如果你的兄弟，一个希伯来男人或希伯来女人，被卖给你，服事你六年；第七年你要让他自由离开你。}在同一本书中，\BibleRef{申 17:15} \BibleQuote{你务必要立上主你的天主所拣选的人为君王统治你；从你的兄弟中，你要立君王统治你；你不可立一个外人，即不是你的兄弟的人，统治你。}又一次，\BibleRef{申 22:1} \BibleQuote{你不可看见你兄弟的牛或羊迷了路，就假装不见；你务必要把它们牵回给你的兄弟。如果你的兄弟离你很远，或者你不认识他，你就把它牵到你家，与你同住，直到你的兄弟来寻找，你再还给他。}宗徒保禄也说，\BibleRef{罗 9:3} \BibleQuote{为了我的兄弟，我的骨肉之亲，就是以色列人，我宁愿自己成为被诅咒的，与基督分离。}此外，那些出于同一家族的人被称为同族的兄弟，就是\OriginalTerm{族系}{\GreekText{πατρία}}，相当于拉丁文的\OriginalTerm{族系}{paternitas}，因为从一个根源产生了众多的后裔。在创世纪中，\BibleRef{创 13:8}我们读到，\BibleQuote{亚巴郎对罗特说：请你不要让我和你之间，以及我的牧人和你的牧人之间发生争吵，因为我们是兄弟。}又一次，\BibleQuote{于是罗特选择了整个约但平原，罗特向东迁移；他们就彼此分开了。}当然，罗特不是亚巴郎的兄弟，而是亚巴郎的兄弟阿兰的儿子。因为特辣黑生了亚巴郎、纳曷尔和阿兰；阿兰生了罗特。我们又读到，\BibleRef{创 12:4} \BibleQuote{亚巴郎离开哈兰时七十五岁。亚巴郎带着妻子撒辣依和侄儿罗特。}但如果你仍然怀疑一个侄子能否被称为儿子，让我给你一个例子。\BibleRef{创 14:14} \BibleQuote{亚巴郎听说他的兄弟被掳，就率领他家生的精练壮丁三百一十八人，直追到丹。}在描述夜袭和杀戮之后，他补充说，\BibleQuote{他把一切财物夺回，也把他的兄弟罗特夺回。}这些足以证明我的主张。但恐怕你会提出一些吹毛求疵的异议，像蛇一样滑脱你的困境，我必须用证明的绳索牢牢捆住你，阻止你嘶嘶作响和抱怨，因为我知道你会想说，你被战胜不是由于圣经真理，而是由于复杂的论证。雅各伯，依撒格和黎贝加的儿子，因害怕他兄弟的诡计而去了美索不达米亚，走近并滚开井口的石头，给他舅舅拉班的羊群饮水。\BibleRef{创 29:11} \BibleQuote{雅各伯吻了辣黑耳，就放声大哭。雅各伯告诉辣黑耳，他是她父亲的外甥，是黎贝加的儿子。}这里是一个已经提到的规则的例子，即一个侄子被称为兄弟。又一次，\BibleRef{创 29:15} \BibleQuote{拉班对雅各伯说：你何必因为你是我的外甥就白白服事我？告诉我，你要什么报酬。}因此，当二十年结束时，雅各伯瞒着岳父，带着妻子和儿子返回本国，拉班在基肋阿得山地追上他，未能找到辣黑耳藏在行李中的神像，雅各伯回答拉班说，\BibleRef{创 31:36} \BibleQuote{我有什么过犯？我有什么罪，你竟这样火热地追赶我？你已经搜遍了我的所有物件，找到了你家的什么物件？把它放在我的兄弟和你的兄弟面前，让他们在我们两人之间评判。}告诉我，当时在场的雅各伯和拉班的兄弟是谁？厄撒乌，雅各伯的兄弟，当然不在那里，而且贝突耳的儿子拉班没有兄弟，尽管他有一个妹妹黎贝加。\par

17. 在圣经中可以找到无数同类实例。但为简洁起见，我愿回到四类弟兄中的最后一类，即那些由亲情而来的弟兄；此类又分为两等：属灵关系的和一般关系的。我之所以说\Emph{属灵}，是因为我们众基督徒都称为弟兄，正如经上所载：\Quote{看，兄弟们同居共处，多么美好，多么快乐！}在另一篇圣咏中，救主说：\Quote{我要向我的弟兄宣扬你的名。}别处又说：\BibleRef{若 20:17} \Quote{你往我弟兄那里去，告诉他们。}我也说\Emph{一般}关系，因为我们同是一位天父的子女，彼此之间有着同样的兄弟纽带。\BibleRef{依 66:5} \Quote{要向那些恨你们的人说：}，先知说，\Quote{你们是我们的弟兄。}宗徒致格林多人书说：\BibleRef{格前 5:11} \Quote{若有称为弟兄的人是淫乱的、或贪婪的、或拜偶像的、或辱骂人的、或酗酒的、或勒索的，这样的人不可和他一同吃饭。}如今我问，福音中主的弟兄应归于哪一类？你们说，他们是按本性为弟兄。但圣经并未如此说；经文既不称他们为玛利亚的儿子，也不称他们为若瑟的儿子。难道说他们是按种族为弟兄？若然如此，那时所有的犹太人都可称为祂的弟兄，这种主张岂不荒谬？他们是因亲密交往和心灵合一的结合而为弟兄吗？倘若这样，谁比那些接受祂私下教诲、被祂称为祂的母亲和弟兄的宗徒们更配称为祂的弟兄？再者，如果所有人作为人都可称为祂的弟兄，那么特意传报\Quote{看，你的弟兄在外边找你}就毫无意义了，因为所有人都同样有权享有这个名号。剩下的唯一选择，是采纳前面的解释，理解为他们被称为弟兄是出于亲属的纽带，而非出于爱和同情，也非种族的特权，更非本性。正如\Person{罗特}被称为\Person{亚巴郎}的弟兄，\Person{雅各伯}被称为\Person{拉班}的弟兄，也如\Person{责罗斐哈得}的女儿们在其弟兄中分得产业，又如\Person{亚巴郎}娶了\Person{撒辣}为妻，称她为妹妹，因为他说：\BibleRef{创 20:11} \Quote{她实在是我的妹妹，同父异母的妹妹}，意即她是他兄弟的女儿，而非姊妹的女儿。否则，我们该如何评论\Person{亚巴郎}——一位义人——娶了同父的女儿呢？圣经在记述古人的历史时，并不以直白语言惊骇我们的耳朵，而宁愿留给读者去推断。而天主后来以法律禁止此行为，并警告说：\BibleRef{肋 18:9} \Quote{谁娶了自己的姊妹，无论是同父的或同母的，并见了她的裸体，这是可耻的事；他应被灭绝。他露了他姊妹的裸体，应负罪债。}\par

18. 有些事物，因你极端的无知，你从未读过，因此你忽略了全部圣经，而将你的疯狂发泄在侮辱童贞女上，如同故事中那个人，他默默无闻，又无法想出善行以获名声，便焚烧了狄安娜神庙；当没有人揭发这种亵渎行为时，据说他亲自四处宣告，他就是纵火者。\Place{厄弗所}的统治者很好奇他为何做这事，他回答说，如果他不能因善行获得名声，那么全人类都应该因恶行记住他。希腊史记载了这件事。但你做得更差。你放火烧了主的身体的圣殿，又污秽了圣神的圣所，你决心从中生出一队四个兄弟和一堆姐妹。总之，你加入犹太人的合唱，说：\Quote{这不是木匠的儿子吗？他的母亲不是叫玛利亚吗？他的兄弟雅各伯、若瑟、西满、犹达不都在我们这里吗？他的姐妹不也全都跟我们在一起吗？如果它们不是一大群，就不会使用\Emph{全部}这个词。}请问，在你出现之前，谁知道这种亵渎？谁认为这理论值两分钱？你得到了你想要的，因罪恶而臭名昭著。至于我，你的对手，虽然我们住在同一城市，用俗话说，我不知道你是白还是黑。我忽略你每本书中充斥的措辞错误。我对你那荒谬的引言只字不提。天哪！我不要求口才，因为你自己没有口才，你向你的兄弟\Person{克拉特里乌斯}求援。我不要求风格的优雅，我只寻求灵魂的纯净：因为对基督徒来说，在言语或行为上引入任何低劣之物，都是最大的语法错误和风格之恶。我的论证即将结束。我将这样对待你，仿佛我至今尚未说服任何人；你会发现你陷入两难困境。显然，我主的兄弟以此名称，正如若瑟被称为他的父亲一样：\BibleRef{路 1:18} \BibleQuote{我同你的父亲忧伤地找寻你。}这是他的母亲说的，不是犹太人。福音作者本人记载，他的父亲和母亲对关于他的话感到惊讶，并且我们已经引用过相似的段落，其中若瑟和玛利亚被称为他的父母。鉴于你愚蠢到说服自己认为希腊文手稿是腐败的，你或许会求助于异文的多样性。因此我来到若望福音，那里清楚写着：\BibleRef{若 1:45} \BibleQuote{斐理伯找到纳塔乃耳，给他说：梅瑟在法律上和先知们所记载的那位，我们找着了，就是若瑟的儿子纳匝肋人耶稣。}你肯定在你的手稿中找到这个。现在告诉我，耶稣怎么可能是若瑟的儿子，明明他是因圣神受生的？若瑟是他真正的父亲吗？你如此迟钝，也不敢那么说。他是他名义上的父亲吗？如果是这样，就让同样的规则应用于被称为兄弟的他们，正如你应用于被称为父亲的若瑟。\par

19. 现在已经避开了暗礁和浅滩，我必须张帆全速到达他的结语。他觉得自己是个一知半解者，在那里提出\Person{戴尔都良}作为证人，并引用\Person{佩塔维的主教维克托利努斯}的话。关于\Person{戴尔都良}，我只需说他不属于教会。至于\Person{佩塔维的主教维克托利努斯}，我断言，从福音书已经证明的事——他所说的主的兄弟并非玛利亚的儿子，而是我所解释的那种兄弟，即由于亲属关系而非本性的兄弟。然而，我们正将精力耗费在琐事上，离开真理的源泉，跟随意见的细小支流。我难道不能为你列举整个古代作者系列吗？\Person{依纳爵}、\Person{波利卡普}、\Person{依勒内}、\Person{殉道者犹斯定}，以及许多其他宗徒的且富有口才的人，他们反对\Person{厄彼翁}、\Person{拜占庭的德奥多托}和\Person{瓦伦提努}，持同样的观点，并写了充满智慧的著作。如果你读过他们的作品，你会更明智。但我认为简短地回应每一点，比继续拖延并将我的书扩展到不合适的长度更好。\par

20. 现在我将攻击那段话，在其中，你想展示你的聪明，你将童贞与婚姻进行了比较。我忍俊不禁，想到了那句谚语：\Emph{你见过骆驼跳舞吗？} \Quote{童贞女更好吗？}你问，\Quote{比已婚的\Person{亚巴郎}、\Person{依撒格}和\Person{雅各伯}更好吗？难道不是每天都有婴孩在天主手中在母亲子宫里成形吗？如果是这样，我们是否应该为玛利亚在分娩后还有丈夫而脸红？如果他们认为这有任何耻辱，那么他们甚至不应该持信心，天主是由童贞女自然分娩而生的。因为按照他们的看法，童贞女通过生殖器官生育天主，比童贞女在分娩后与自己的丈夫结合更为可耻。}如果你愿意，再加上\Person{赫尔维狄乌斯}，其他的本性屈辱：子宫九个月变大，疾病，分娩，血，襁褓。想象婴儿在包裹的胎膜中。在你的画面中引入硬冷的马槽，婴儿的哭喊，第八天的割损礼，取洁的时期，这样他可以被证明是不洁的。我们不脸红，我们不沉默。他为我受的屈辱越大，我欠他的越多。当你给出所有细节时，你将能提出的没有比十字架更可耻的，我们承认它，我们相信它，并籍着它克胜我们的仇敌。\par

21. 但我们既然不否认所记载的，同样也拒绝所没有记载的。我们相信天主由童贞女所生，因为我们读到这事。至于玛利亚在生产之后结了婚，我们不相信，因为我们没有读到这事。我们这样说并非要谴责婚姻，因为童贞本身正是婚姻的果实；而是因为在论及圣人时，我们不可妄断。若我们以可能性作为判断的标准，我们或许可以主张若瑟有好几位妻子，因为亚巴郎有，雅各伯也有，并且主的弟兄们是这些妻子所生的——这是一种某些人以鲁莽（源于大胆而非虔敬）所持的虚构。你说玛利亚没有保持童贞；我更进一步主张：若瑟本人因玛利亚也是童贞的，如此从童贞的婚姻中生下了一个童贞之子。因为若一个圣洁的人不被视为犯奸淫，且没有任何地方记载他另娶了妻子，他乃是玛利亚的监护人——人们以为他是她的丈夫，而实际上并非如此——那么结论是：那被认为配称为主之父的人，一直保持童贞。\par

22. 现在我要在童贞与婚姻之间作一比较，恳请读者不要以为我在赞美童贞时稍稍贬低了婚姻，或把旧约的圣人与新约的圣人分隔开来——也就是说，把那些有妻子的人与那些完全禁绝女性怀抱的人分开；我反倒认为，按照时代与环境的不同，前者适用一种规则，而我们这些万世终结已临到的人则适用另一种规则。只要那法律仍在：\BibleRef{创 1:28}\Quote{你们要生育繁殖，充满大地}；以及\Quote{以色列中不生育的妇女是有祸的，因为她不在以色列中留下后裔}，他们全都男婚女嫁，离开父母，成为一体。但一旦雷声般的话被听到：\BibleRef{格前 7:29}\Quote{时间是缩短的了，今后有妻子的要像没有一样}；我们紧紧依附于主，与祂成为一神。为什么呢？因为\Quote{没有娶妻的，所挂虑的是主的事，想怎样悦乐主；娶了妻的，所挂虑的是世俗的事，想怎样悦乐妻子。有妻子的妇女和童贞女也有区别。没有出嫁的，挂虑主的事，一心使身体和灵魂保持圣洁；出嫁的，挂虑世俗的事，想怎样悦乐丈夫。}你为何吹毛求疵？为何抗拒？这是蒙选之器所说的；他告诉我们妻子与童贞女之间是有区别的。请注意那状态该是何等幸福，连性别的区分都消失了。童贞女不再被称为女人。\BibleRef{格前 7:34}\Quote{没有出嫁的，挂虑主的事，一心使身体和灵魂保持圣洁。}童贞女被定义为在身体和灵魂上都圣洁的人，因为若一个女人在心思中结了婚，仅有童贞的肉身是没有益处的。\par

\Quote{但出嫁的，挂虑世俗的事，想怎样悦乐丈夫。}你以为整日祈祷禁食的人，与必须在她丈夫到来时化妆、扭捏作态、假装亲热的人之间没有区别吗？童贞女的目标是显得不那么标致；她会损害自己，以掩盖天生的姿色。已婚妇女对着镜子涂脂抹粉，侮辱她的造物主，试图获取超乎自然的美。然后便是婴孩的呀呀学语、喧闹的家室、孩子等待她的话、等待她的亲吻、算计开销、准备应付开支。一方面你会看见一群厨子，束带备战，猛攻肉食；那里你可以听见众多织工的嗡嗡声。与此同时，消息传来，丈夫和他的朋友们到了。妻子像燕子一样飞遍整个屋子。\Quote{她得照料一切。沙发平滑吗？地板扫干净了吗？花瓶里有花吗？晚餐准备好了吗？}请问，这一切之中哪里有思念天主的余地？这些是幸福的家吗？在有鼓声、笛琴的喧闹、钹的铿锵之处，还能找到任何对天主的敬畏吗？食客被斥责，却以此为荣。接着进来半裸的情欲受害者，成为每个淫荡眼光的靶子。不幸的妻子要么喜欢她们而丧亡，要么不喜欢而激怒丈夫。由此产生不睦——离婚的温床。或者你找一个没有这些事的家庭给我看，那真是\Emph{\OriginalTerm{稀罕之物}{rara avis}}！然而即便如此，家务管理、子女教育、丈夫的需要、仆人的管教，都不可避免地会使人分心，不能思念天主。\BibleRef{创 18:11}\Quote{撒辣已过了月经的年龄}——经上这样说——之后亚巴郎受命：\BibleRef{创 21:12}\Quote{凡撒辣对你说的，你都要听从她的话。}那不受怀孕和生产的焦虑折磨，且已过了更年期、不再履行女性职能的女子，便脱离了天主的咒诅；她的欲望不再指向丈夫，相反，丈夫反成了她的属下，上主的声音命令他：\Quote{凡撒辣对你说的，你都要听从她的话。}于是他们开始有时间祈祷。因为只要婚姻的债还在偿还，热切的祈祷就会被忽视。\par

23. 我不否认，在寡妇和已婚妇女中都可找到圣善的妇女；但她们或是已不再为妻子者，或是在婚姻的紧密联系中仍效法童贞之洁德者。宗徒——基督在他内说话——曾简要地为此作证，他说：\BibleRef{格前 7:34} \BibleQuote{没有丈夫的妇女，关心主的事，为叫身心圣洁；至于已出嫁的，关心世俗的事，想怎样悦乐丈夫。} 他在这事上任由我们运用理性。他不将必要性加于任何人，也不引人陷入圈套；他只劝勉那适宜的事，当他愿意众人像他自己一样时。确实，他从主那里没有领受关于童贞的诫命，因为那恩宠超越了人本身的能力；而强迫人违背本性，换句话说，要你们成为天使那样，那会带有不谦逊的意味。正是这种天使般的纯洁确保了童贞的最高赏报；而宗徒本可轻视一种不涉及罪过的生活方式。然而，在紧接着的上下文中，他加上：\BibleRef{格前 7:25} \BibleQuote{论到童身的人，我没有主的命令，但我既蒙主仁慈，作为一个忠信的人，说出我的意见：所以，为了现时的困境，我以为为人这样倒好。} \Emph{现时的困境}是什么意思？\BibleQuote{怀孕的和乳养孩子的有祸了！} 树木之所以生长，是为了被砍伐；田地之所以播种，是为了被收割。世界已经满了，人口对于土地来说太多了。我们每天都因战争而被砍伐，被疾病夺去，被海难吞没，尽管我们为财产的界限而彼此诉讼。只有那些跟随羔羊、没有玷污自己衣裳的人，才是在普遍规律之外的一个例外，因为他们保持了童贞状态。请注意\Emph{玷污}的含义。我不敢解释它，以免\Person{厄耳味丢}辱骂。我同意你所说的，有些童贞女不过是酒馆女郎；我甚至说，其中也可找到奸妇；你无疑会更惊讶地听到，有些神职人员是旅店老板，有些隐修士不贞洁。谁不会立刻明白，酒馆女郎不可能是童贞女，奸夫不可能是隐修士，神职人员不可能是旅店老板？如果假冒的童贞有缺陷，我们该责怪童贞本身吗？就我而言，略过其他人而回到童贞女，我主张：那从事叫卖的女子，虽然就我所知可能在身体上是童贞，但在精神上已不再是童贞。\par

24. 我变得好辩，像讲坛上的演说家一样稍事娱乐。是你强迫了我，\Person{厄耳味丢}；因为尽管今日福音大放光明，你却主张童贞与婚姻状态具有同等的荣耀。并且因为我认为，你发现真理对你来说过于强大，你就会转而贬低我的生活和毁谤我的品格（懦弱的妇人被主人压制时，习惯在角落闲言碎语），我将先发制人。我向你保证，我会把你的辱骂视为极大的荣誉，因为攻击我的同一张嘴唇曾贬抑玛利亚，而我，主的仆人，荣幸地享有与祂母亲同样的犬吠般的雄辩。\par

\SourceRecord{Translated by W.H. Fremantle, G. Lewis and W.G. Martley.；Nicene and Post-Nicene Fathers, Second Series；Vol. 6.；Edited by Philip Schaff and Henry Wace.；Buffalo, NY: Christian Literature Publishing Co.,；1893.；<http://www.newadvent.org/fathers/3007.htm>.}
